Lista de verificación para usar antes de enviar una solución y cuando el veredicto no es \textit{Accepted}.

\textbf{Antes de enviar.}
\begin{props}
  \item Escribir casos de prueba propios si el ejemplo no alcanza. Si el límite de tiempo está justo, generar el caso máximo y medir.
  \item Quitar la salida de depuración. Para depurar conviene escribir a \texttt{sys.stderr}, que no afecta la salida.
  \item Confirmar que se envía el archivo correcto y que el formato de salida coincide (espacios, saltos de línea, mayúsculas como \texttt{YES}/\texttt{Yes}).
\end{props}

\textbf{WA (Wrong Answer).}
\begin{props}
  \item Releer el enunciado completo. ¿Se entendió bien el problema? ¿El algoritmo funciona para todo el rango de entrada y para los casos especiales (\(n = 1\), lista vacía, todos iguales, negativos)?
  \item Escribir una solución lenta pero correcta (fuerza bruta) y compararla con la rápida sobre casos aleatorios pequeños.
  \item ¿Se limpian las estructuras globales entre casos de prueba? ¿Hay variables sin inicializar o confundidas (\(n\)/\(m\), \(i\)/\(j\))?
  \item \texttt{/} devuelve \texttt{float} y pierde precisión con enteros grandes: usar \texttt{//} cuando el resultado debe ser entero. Comparar \texttt{float} con tolerancia, no con \texttt{==}.
  \item Matrices: \texttt{[[0] * m] * n} repite \textbf{la misma fila} \(n\) veces. Usar \texttt{[[0] * m for \_ in range(n)]}.
  \item \texttt{b = a} no copia la lista: usar \texttt{a[:]}, \texttt{list(a)} o \texttt{copy.deepcopy}. Un argumento por defecto mutable (\texttt{def f(x, memo=[])}) se comparte entre llamadas.
  \item Los índices negativos no dan error (\texttt{a[-1]} es el último elemento), por lo que un error de borde puede pasar sin avisar.
  \item Revisar \texttt{return} vs \texttt{continue} vs \texttt{break}, y \texttt{<} vs \texttt{<=}. No modificar una lista mientras se itera sobre ella.
  \item Agregar \texttt{assert} y volver a enviar. Explicar el algoritmo a un compañero o pedirle que lea el código.
\end{props}

\textbf{TLE (Time Limit Exceeded).}
\begin{props}
  \item ¿Hay un posible bucle infinito? ¿Cuál es la complejidad real? Un límite de tiempo alto no implica que no se espere algo eficiente (ver la tabla de complejidad esperada).
  \item Usar entrada y salida rápidas (\texttt{sys.stdin}, un solo \texttt{print} con \texttt{join}).
  \item Evitar operaciones ocultas de costo \(O(n)\): \texttt{x in lista} (usar \texttt{set} o \texttt{dict}), \texttt{lista.pop(0)} o \texttt{insert(0, x)} (usar \texttt{deque}), \texttt{s += c} en un ciclo sobre strings (acumular en una lista y usar \texttt{join}), y cortar listas (\texttt{a[i:]}) en cada llamada recursiva (pasar índices).
  \item Si la plataforma lo permite, enviar con PyPy.
\end{props}

\textbf{RTE (Runtime Error).}
\begin{props}
  \item \texttt{IndexError}: lectura fuera de rango. \texttt{KeyError}: clave ausente en un \texttt{dict} (considerar \texttt{defaultdict} o \texttt{.get}). \texttt{ZeroDivisionError}: división o módulo por 0.
  \item \texttt{RecursionError}: subir \texttt{sys.setrecursionlimit} o pasar a una pila explícita (ver Recursión profunda).
  \item \texttt{ValueError}: \texttt{int(\textquotesingle\textquotesingle)} por una línea vacía, o conversión de un entero de más de 4300 dígitos a texto (ver Recomendaciones generales). \texttt{EOFError}: se intentó leer más datos de los que hay.
\end{props}

\textbf{MLE (Memory Limit Exceeded).}
\begin{props}
  \item Estimar la memoria máxima necesaria. Una lista de \(10^6\) elementos ocupa unos 8 MB solo en referencias, y cada entero distinto ocupa además unos 28 bytes. Para volúmenes grandes conviene \texttt{array} (por ejemplo, \texttt{array('i')} usa 4 bytes por elemento) o \texttt{bytearray}.
  \item Limpiar o reutilizar las estructuras entre casos de prueba. Usar generadores en vez de construir listas intermedias cuando solo se recorren una vez.
\end{props}
